\section{本章小结}
\label{sec:rkhs-summary}

章首的问题是：有限个点值怎样约束一条函数。本章的回答从候选空间开始。有限特征空间把插值化为设计矩阵方程，但维数一变，可行性与唯一性随之改变，零残差也不保证测点之外的误差为零。要在无限维空间中选函数，候选空间必须让点值有确定意义且受范数控制；$L^2$两项都不满足，$H^1(0,1)$与$H_{\min}$都满足。

把每个点评估都连续的Hilbert函数空间称为RKHS之后，Riesz表示把读数写成与核截面的内积$f(x)=\langle K_x,f\rangle$，核截面的内积就是再生核。再生核共轭对称、正定，对角值$K(x,x)$给出点评估的范数，核截面的组合在空间中稠密。反过来，任何正定核经过半内积、零范数商与完备化，唯一确定一个以它为再生核的RKHS；补进来的极限元素由点值估计识别为函数。新核可以由特征映射、非负线性组合与逐点乘积得到，同一个核可以有维数不同的特征表示，RKHS的范数等于能给出该函数的最短参数长度。

有限观测问题由观测映射$S$与Gram矩阵$\bm G=SS^*$完整描述。观测零空间$N=\ker S$与核截面张成的$M$互为正交补，投影到$M$保留读数而不增大范数。数据可插值当且仅当$\bm y$属于$\bm G$的像；可行时，最小范数插值$f^+=\sum_jc_jK_{x_j}$唯一，即使$\bm G$奇异、系数不唯一也是如此。没有插值函数与插值函数不唯一由不同对象决定，前者看$\bm G$的像，后者看$N$。

表示误差的比较必须说明在哪个空间、用哪些观测、按哪个范数。在RKHS内，子空间$V$的投影误差乘以$\sqrt{K(x,x)-K_V(x,x)}$控制逐点误差；核截面的坐标由点值读出，限定点集的截面是否稠密取决于核的连续性与点集是否足以排除非零函数。一般有限特征可以更经济，但其欧氏内积给出的有限特征核不是原核，原空间中子空间的核要按$H$的Gram矩阵加权。

加入全支撑的有限测度后，紧度量空间上连续正定核的积分算子是紧、自伴、正算子，Mercer定理给出核的绝对一致展开$K(x,y)=\sum_j\lambda_j\phi_j(x)\overline{\phi_j(y)}$。$\{\sqrt{\lambda_j}\phi_j\}$是RKHS的标准正交基，同一组谱坐标在$L^2$与RKHS中的权重相差$1/\lambda_j$。截断前$N$个谱方向时，RKHS单位球上的$L^2$平方误差统一不超过$\lambda_{N+1}$，RKHS误差却没有统一趋零的界。读数带噪声时的损失与正则化选择属于卷五，随机观测的统计误差属于卷八；它们使用的正是本章的正交分解与有限表示。
